\documentclass[11pt]{article}
\usepackage[letterpaper,margin=1in]{geometry}
\usepackage{times}
\pagestyle{empty}
\begin{document}
\noindent 27 November 2026
\vspace{1.5em}

\noindent The Editors\\
\textit{Journal of Biomechanics}
\vspace{1.5em}

\noindent Dear Editors,
\vspace{1em}

\noindent I am submitting the attached manuscript, ``Device-normalised geometry for
muscle architecture measurement in B-mode ultrasound,'' for consideration in the
\textit{Journal of Biomechanics}.
\vspace{1em}

\noindent Measuring pennation angle, fascicle length and muscle thickness from B-mode
ultrasound is well covered by existing tools, but those tools assume a known
setup: a single device, a consistent export, and a scale the operator supplies.
This paper handles the case where none of that is available. The input is an
unlabelled screenshot from one of six different ultrasound systems, with no
metadata and no stated scale.
\vspace{1em}

\noindent Two findings may be of interest to your readers. The first is that two of the six
systems export a horizontally resampled image, so the pixels are not square and
any angle measured in pixel coordinates is biased by a fixed factor. That factor
can be computed from the published transducer aperture, so the correction needs
no labelled data. The second is that the conventional fascicle-length estimator,
which extends a detected fragment until it meets both aponeuroses, carries a bias
of $1/(1-\cot^{2}\gamma\tan^{2}h)$ that depends on how much the aponeuroses
converge. Human raters do not reproduce that bias. Measuring against the bisector
of the two aponeuroses removes the dependence, which is what allows one
calibration constant to carry from a static benchmark to a dynamic sequence of a
different muscle without refitting.
\vspace{1em}

\noindent On a 35-image benchmark measured independently by seven expert raters, the method
scores 0.196 under the UMUD error formula. DL\_Track\_US scores 0.331 on the same
images, SMA scores 0.409, and an individual expert rater scores 0.264 by
leave-one-out. All automated methods are scored against the same consensus on the
same images, using measurements published by the benchmark's own authors rather
than re-implementations of my own.
\vspace{1em}

\noindent This manuscript is original, has not been published before, and is not under
consideration anywhere else. I am the sole author. I have no conflicts of
interest and the work had no external funding. The evaluation code, model weights
and split manifests are publicly available.
\vspace{1em}

\noindent Thank you for considering it.
\vspace{2em}

\noindent Sincerely,
\vspace{1.5em}

\noindent \textbf{Colton Downing}\\
University of Southern California\\
Los Angeles, CA, USA\\
cjdownin@usc.edu\\
ORCID: 0009-0009-8874-1914
\end{document}
